\chapter*{Sources and further reading}
\addcontentsline{toc}{chapter}{Sources and further reading}
Vendor descriptions were inspected for substrate specificity, directionality, and activity qualifications, not for the assigned rates. The ligase paper's abstract and introduction were inspected for the adenylation mechanism. The cited IUBMB sections were inspected for kinetic definitions and transient versus steady-state rates. Access checked 20 September 2026.

\begingroup\let\chapter\section
\begin{thebibliography}{9}
\bibitem{ecori} New England Biolabs. \emph{EcoRI}, R0101: recognition site, cleavage position, and specificity qualifications.
\url{https://www.neb.com/en-us/products/r0101-ecori}.
\bibitem{ligase} New England Biolabs. \emph{T4 DNA Ligase}, M0202: substrate and cofactor description.
\url{https://www.neb.com/en/products/m0202-t4-dna-ligase}.
\bibitem{ligasemechanism} K. Shi et al. \emph{T4 DNA ligase structure reveals a prototypical ATP-dependent ligase with a unique mode of sliding clamp interaction}. Nucleic Acids Research 46(19), 10474--10488 (2018).
\url{https://doi.org/10.1093/nar/gky776}.
\bibitem{polymerase} New England Biolabs. \emph{T4 DNA Polymerase}, M0203: direction, primer requirement, and exonuclease activity.
\url{https://www.neb.com/products/m0203-t4-dna-polymerase}.
\bibitem{kinetics} IUBMB. \emph{Symbolism and Terminology in Enzyme Kinetics} (1981), sections 4.1--4.2.
\url{https://iubmb.qmul.ac.uk/kinetics/ek4t6.html}.
\bibitem{transient} IUBMB. \emph{Symbolism and Terminology in Enzyme Kinetics} (1981), introductory definitions of steady-state and transient behavior.
\url{https://iubmb.qmul.ac.uk/kinetics/ek1t3.html}.
\end{thebibliography}
\endgroup
